\documentclass[10pt,a4paper]{amsart}
\usepackage[margin=19mm]{geometry}
\usepackage{amsmath,amssymb,mathtools}
\usepackage[scr=rsfs,cal=euler]{mathalfa}
\usepackage{xfrac}
\usepackage[unicode,hidelinks,bookmarks=false]{hyperref}
\newcommand{\ie}{i.e.\ }
\newcommand{\cf}{cf.\ }
\newcommand{\resp}{respectively\ }
\newcommand{\Subsection}{\subsection}
\providecommand{\nobreakdash}{\nobreak\hbox{-}}
\setlength{\emergencystretch}{2em}
\multlinegap0pt
\usepackage{xcolor}
\usepackage{amssymb}
\usepackage[normalem]{ulem}
\usepackage[shortlabels]{enumitem}
\usepackage{float}
\usepackage{dsfont}
\usepackage{csquotes}
\usepackage{algorithm}\usepackage{algpseudocode}
\usepackage{mathtools}
\newtheorem{lem}{Lemma}
\title{Mixed-identity-freeness and primitivity of group rings}
\author{Felipe I. Flores}
\date{Source: arXiv:2607.28316v1 (CC BY 4.0)}
\begin{document}
\maketitle

\noindent\emph{Source and licence.} Mixed-identity-freeness and primitivity of group rings. Version arXiv:2607.28316v1 (CC BY 4.0), distributed by arXiv under the Creative Commons Attribution 4.0 International licence, http://creativecommons.org/licenses/by/4.0/, as recorded in the arXiv licence metadata for this exact version and shown on the abstract page https://arxiv.org/abs/2607.28316v1. Full licence notice: \texttt{licenses/arxiv-2607.28316-CC-BY-4.0.txt}.\par\smallskip

\noindent\textit{Editorial scope.} Counted unit: the article's lem lem (source0.tex lines 366--368). The article's complete original proof of it is delivered with it (source0.tex lines 370--398, one \texttt{proof} environment of 29 lines). No mathematics is written, repaired or re-derived: the statement and the proof are the article's own lines, in the article's own notation, and no retained line is substituted or re-flowed.  No further source-local context is needed: every cross-reference of every retained line resolves inside the two delivered spans.  Nothing was omitted from any retained span, so the package carries no omission record.  No retained line contains a citation, so the wrapper carries no bibliography and no bibliography entry.  No retained line contains a figure environment, an image-inclusion command or drawing code, so no image is delivered and the package declares no asset dependency.  Editorial formatting only, in addition: an AMS \texttt{amsart} preamble that restates the article's own declarations and macros used by the retained lines, namely the environments \texttt{lem} on separate counters, together with the standard packages those lines need.  The retained environments print in this document as Lemma 1 in order of appearance, whereas the article prints the same items with the numbering of its own full document.  The complete native source of the article as distributed in the arXiv e-print for this exact version is bundled unchanged as \texttt{sources/206406/source0.tex}.
\begin{lem}
    Let $G$ be mixed-identity-free. Given finitely many nontrivial elements $w_1,\ldots,w_m\in G*\mathbb F_n$, there exists $g_1,\ldots,g_n\in G$ such that $w_j(g_1,\ldots,g_n)\not=1$, for all $1\leq j\leq m$.
\end{lem}

\begin{proof}
Introduce new free variables $y_2,\ldots,y_m$. Define recursively
$$
W_1=w_1\quad\text{ and }\quad W_j=[W_{j-1},y_j^{-1}w_jy_j],
$$
where $2\leq j\leq m$.

We claim that each $W_j$ is nontrivial. Suppose inductively that
$W_{j-1}\neq1$ and put
$$
H=G*F(t_1,\ldots,t_n,y_2,\ldots,y_{j-1}).
$$
Then $W_{j-1},w_j$ are nontrivial elements of $H$. In the free product $H*\langle y_j\rangle$, we have
$$
W_j
=
W_{j-1}^{-1}y_j^{-1}w_j^{-1}y_j
W_{j-1}y_j^{-1}w_jy_j.
$$
This is a reduced free-product normal form: all displayed $H$-syllables are
nontrivial and alternate with nontrivial powers of $y_j$. Hence
$W_j\neq1$.

Because $G$ is MIF, there are $g_1,\ldots,g_n,h_2,\ldots,h_m\in G$ such that
$$
W_m(g_1,\ldots,g_n,h_2,\ldots,h_m)\neq1.
$$
If $w_j(g_1,\ldots,g_n)=1$ for some $j$, then $W_j$ evaluates to $1$, and consequently, all subsequent words $W_k$ evaluate to $1$, including $W_m$. This is a contradiction. Therefore, every $w_j(g_1,\ldots,g_n)\neq1$.
\end{proof}

\end{document}
